% =====================================================================
%  Lecture 26 -- Future 6G Technologies                   (Persian edition)
% =====================================================================
\documentclass[11pt,a4paper]{article}
\providecommand\ACNroot{../..}
\input{\ACNroot/templates/lecture-fa}

\ACNset{number}{26}
\ACNset{title}{فناوری‌های آیندهٔ نسل ششم}
\ACNset{shorttitle}{فناوری‌های نسل ششم}
\ACNset{titleother}{Future 6G Technologies}
\ACNset{part}{بخش پنجم -- صفحه‌ها و چشم‌انداز (درس‌های ۲۳ تا ۲۶)}
\ACNset{time}{۶ تا ۸ ساعت (۳ ساعت مطالعه، ۳ تا ۵ ساعت تمرین)}
\ACNset{level}{پیشرفته}
\ACNset{prereq}{درس‌های ۱۴، ۲۳ و ۲۵؛ ظرفیت شانون و انتشار در فضای آزاد}
\ACNset{gen}{6}

\begin{document}
\ACNcover

% ---------------------------------------------------------------------
\section{اهداف یادگیری}
پس از مطالعهٔ این درس باید بتوانید:
\begin{objectives}
  \item چارچوب \lr{IMT-2030} یعنی سناریوهای کاربرد و جنبه‌های فراگیر آن را \textbf{شرح دهید}.
  \item فناوری‌های نامزد اصلی را \textbf{توضیح دهید}: طیف زیرتراهرتز، \acro{MIMO} انبوه بدون سلول، سطوح هوشمند بازپیکربندی‌پذیر، حسگری و ارتباط یکپارچه،
        شبکه‌های غیرزمینی و طراحی هوش‌مصنوعی‌بنیاد.
  \item افت مسیر فضای آزاد در بسامدهای بالا و بهرهٔ ظرفیت پهنای باند بیشتر را تحت محدودیت توان \textbf{محاسبه کنید}.
  \item بهرهٔ \acro{RIS}، تفکیک‌پذیری حسگری و تأخیر انتشار ماهواره را \textbf{برآورد کنید}.
  \item \textbf{ارزیابی کنید} کدام ادعاها دربارهٔ نسل ششم تثبیت‌شده‌اند و کدام هنوز پژوهشی‌اند.
\end{objectives}

% ---------------------------------------------------------------------
\section{دانش پیش‌نیاز}
به الزامات نسل پنجم و طراحی \acro{NR} (درس ۱۴)، بودجه‌های صفحهٔ کاربر درس ۲۳ و روندهای تکامل درس ۲۵ نیاز دارید، و نیز به فرمول ظرفیت شانون و معادلهٔ
انتشار فضای آزاد (فریس).

\paragraph{واژه‌های کلیدی} واژگان این درس همراه با معادل انگلیسی آن‌ها:
\begin{keyterms}
\lr{IMT-2030} & \lr{IMT-2030} & چارچوب \acro{ITU} برای نسل ششم (توصیه‌نامهٔ \lr{M.2160}). \\
زیرتراهرتز & \lr{Sub-THz} & بسامدهای حامل تقریباً ۱۰۰ تا ۳۰۰ گیگاهرتز. \\
\acro{MIMO} انبوه بدون سلول & \lr{Cell-free massive MIMO} & نقاط دسترسی توزیع‌شدهٔ بسیار که با هم به همهٔ کاربران خدمت می‌دهند. \\
سطح هوشمند بازپیکربندی‌پذیر & \lr{RIS} & سطحی از عناصر غیرفعال با فاز کنترل‌پذیر که بازتاب‌ها را شکل می‌دهد. \\
حسگری و ارتباط یکپارچه & \lr{ISAC} & استفاده از همان سیگنال‌ها و سخت‌افزار برای حسگری راداری و ارتباط. \\
شبکهٔ غیرزمینی & \lr{NTN} & ماهواره‌ها و سکوهای ارتفاع بالا به‌عنوان بخشی از شبکهٔ موبایل. \\
هوش‌مصنوعی‌بنیاد & \lr{AI-native} & طراحی‌شده با اجزای یادگیرنده و خطوط لولهٔ داده از ابتدا. \\
تفکیک‌پذیری فاصله & \lr{Range resolution} & کوچک‌ترین اختلاف فاصله میان دو هدف قابل تمایز. \\
رژیم محدود به توان & \lr{Power-limited regime} & جایی که پهنای باند بیشتر دیگر ظرفیت را چندان افزایش نمی‌دهد. \\
انرژی به ازای بیت & \lr{Energy per bit} & انرژی مصرفی تقسیم بر بیت‌های تحویل‌شده (\lr{J/bit}). \\
\end{keyterms}

% ---------------------------------------------------------------------
\section{مقدمه و انگیزه}
هر نسل حدود ده سال پس از نسل قبلی تعریف شده است و نسل ششم همین آهنگ را دنبال می‌کند: \acro{ITU} چارچوب \lr{IMT-2030} را در ۲۰۲۳ منتشر کرد
\cite{itu2023m2160} و \lr{3GPP} کار مطالعاتی نسل ششم را در انتشار ۲۰ آغاز کرد و مشخصات هنجاری در انتشارهای بعدی انتظار می‌رود. همان‌طور که درس ۲۵
نشان داد، گام‌های بعدی ظرفیت به پهنای باندهای بسیار گسترده‌تر نیاز دارند، هزینه به ازای بیت باید همچنان کاهش یابد و بخش‌های بزرگی از جهان هنوز
بدون اتصال‌اند. هم‌زمان از شبکه انتظار می‌رود فراتر از جابه‌جایی بیت‌ها عمل کند: محیط خود را حس کند، میزبان هوش مصنوعی باشد و پایدار و
تاب‌آور باشد \cite{saad2020vision}.

بخش بزرگی از نسل ششم هنوز پژوهشی است. بنابراین این درس چارچوب توافق‌شده را از فناوری‌های نامزد جدا می‌کند و مدل‌های فیزیکی ساده‌ای ارائه
می‌دهد که به شما اجازه می‌دهند ادعاهای مربوط به آن‌ها را خودتان داوری کنید.

% ---------------------------------------------------------------------
\section{مفاهیم بنیادی}

\subsection{چارچوب \lr{IMT-2030}}
توصیه‌نامهٔ \lr{ITU-R M.2160} \cite{itu2023m2160} شش سناریوی کاربرد تعریف می‌کند (\cref{fig:usage}). سه سناریو مثلث نسل پنجم را گسترش می‌دهند:
\emph{ارتباط فراگیر} (از \acro{eMBB})، \emph{ارتباط فوق‌قابل‌اطمینان و کم‌تأخیر} (\acro{HRLLC}، از \acro{URLLC}) و \emph{ارتباط انبوه} (از \acro{mMTC}). سه
سناریو جدیدند: \emph{اتصال همه‌جایی}، \emph{هوش مصنوعی و ارتباط} و \emph{حسگری و ارتباط یکپارچه}. چهار جنبهٔ فراگیر بر همه حاکم‌اند: پایداری، اتصال
بی‌اتصال‌ها، هوش فراگیر، و امنیت و تاب‌آوری. الزامات کمّی در گام بعدی \acro{ITU} تعریف می‌شوند، پس اعدادی مانند «۱ ترابیت بر ثانیه» که در منابع
پژوهشی آمده‌اند اهداف پژوهشی‌اند، نه الزام.

\begin{acnfigure}
  \resizebox{\linewidth}{!}{\acnfig{lecture-26/usage-scenarios}}
  \caption{سناریوهای کاربرد و جنبه‌های فراگیر \lr{IMT-2030} (بر اساس \lr{ITU-R M.2160}).}
  \label{fig:usage}
\end{acnfigure}

\subsection{فناوری‌های نامزد}
\begin{itemize}
  \item \textbf{طیف جدید}: باند میانی بالایی (حدود ۷ تا ۲۴ گیگاهرتز) برای ظرفیت گسترده و باندهای زیرتراهرتز برای نرخ‌های بسیار بالا در برد کوتاه
        \cite{rappaport2019above100}.
  \item \textbf{\acro{MIMO} انبوه بدون سلول}: نقاط دسترسی بسیار که بدون مرز سلولی به‌صورت همدوس به کاربران خدمت می‌دهند و خدمت یکنواخت فراهم می‌کنند
        \cite{ngo2017cellfree}.
  \item \textbf{سطوح هوشمند بازپیکربندی‌پذیر}: بازتابنده‌های کنترل‌پذیر که با توان کم مسیرهایی پیرامون موانع می‌سازند \cite{wu2019irs}.
  \item \textbf{حسگری و ارتباط یکپارچه}: رادیو به‌عنوان رادار برای مکان‌یابی، نقشه‌برداری و تشخیص حرکت یا اشیا \cite{liu2022isac}.
  \item \textbf{شبکه‌های غیرزمینی}: ماهواره‌های \acro{LEO} و سکوهای ارتفاع بالا یکپارچه با سامانهٔ زمینی \cite{tr38811}.
  \item \textbf{رابط هوایی و شبکهٔ هوش‌مصنوعی‌بنیاد}: گیرنده‌های یادگیرنده، مدیریت پرتو و بهره‌برداری خودکار.
\end{itemize}

% ---------------------------------------------------------------------
\section{توضیحات فنی تفصیلی}

\subsection{ارتباط زیرتراهرتز}
بالای ۱۰۰ گیگاهرتز پهنای باندهای پیوستهٔ ده‌ها گیگاهرتزی در دسترس است، اما برای آنتن‌هایی با بهرهٔ ثابت افت مسیر با بسامد زیاد می‌شود، جذب مولکولی
در بسامدهای خاصی پدید می‌آید و انسداد توسط دیوارها، بدن‌ها و حتی شاخ و برگ شدید است. آرایه‌های بسیار بزرگ با پرتوهای باریک افت را جبران می‌کنند
(\cref{sec:fspl})، و این جست‌وجو و ردیابی پرتو را به مسئلهٔ اصلی تبدیل می‌کند. سخت‌افزار نیز دشوار است: تقویت‌کننده‌های توان بازدهی کمی دارند،
نویز فاز با بسامد افزایش می‌یابد و نمونه‌برداری ده‌ها گیگاهرتز به مبدل‌های کم‌دقت نیاز دارد.

\subsection{شبکه‌های بدون سلول}
در یک شبکهٔ سلولی، کاربران لبهٔ سلول از تداخل و سیگنال ضعیف رنج می‌برند. در یک شبکهٔ بدون سلول، یک واحد پردازش مرکزی نقاط دسترسی توزیع‌شدهٔ
بسیاری را هماهنگ می‌کند که همه به همهٔ کاربران خدمت می‌دهند، با پیش‌کدگذاری محلی و فرانت‌هال به واحد مرکزی. نتیجه نرخ‌های یکنواخت‌تر است؛ هزینه‌ها
فرانت‌هال و هم‌گام‌سازی‌اند -- گسترشی از ایده‌های تقسیم \acro{CU}/\acro{DU} و هماهنگی در درس ۲۳.

\subsection{سطوح هوشمند بازپیکربندی‌پذیر}
یک \acro{RIS} صفحه‌ای از عناصر زیرطول‌موجی بسیار است که فاز بازتاب آن‌ها به‌صورت الکترونیکی تنظیم می‌شود. با پیکربندی درست، یک مسیر بازتابی قوی
می‌افزاید، مثلاً پیرامون یک ساختمان مسدودکننده. توان دریافتی آن با مجذور تعداد عناصر رشد می‌کند (\cref{sec:ris})، اما از حاصل‌ضرب دو افت مسیر (فرستنده
تا سطح، سطح تا گیرنده) رنج می‌برد، پس بیشترین کمک را نزدیک فرستنده یا گیرنده دارد.

\subsection{حسگری و ارتباط یکپارچه}
پهنای باندهای گسترده و آرایه‌های بزرگ دقیقاً همان چیزی است که رادارها نیاز دارند. \acro{ISAC} از شکل‌موج‌های ارتباطی یا سخت‌افزار مشترک برای اندازه‌گیری
فاصله، زاویه و سرعت اشیا استفاده می‌کند. کاربردها شامل مکان‌یابی بدون زیرساخت اختصاصی، پایش ترافیک و نقشه‌برداری محیط است که خود به مدیریت پرتو
کمک می‌کند. این فناوری پرسش‌هایی دربارهٔ حریم خصوصی نیز پیش می‌کشد، چون شبکه می‌تواند افرادی را که هیچ دستگاهی همراه ندارند حس کند.

\subsection{شبکه‌های غیرزمینی}
انتشار ۱۷ \acro{NR} روی ماهواره (\acro{NTN}) را با سازگاری‌هایی برای تأخیرهای بزرگ و جابه‌جایی داپلر معرفی کرد. منظومه‌های \acro{LEO} پوشش جهانی با
تأخیر رفت‌وبرگشت چند ده میلی‌ثانیه‌ای فراهم می‌کنند (\cref{sec:leo})؛ با یکپارچگی با شبکه‌های زمینی، به «اتصال بی‌اتصال‌ها» و تاب‌آوری پس از بلایا
می‌پردازند.

\subsection{شبکه‌های هوش‌مصنوعی‌بنیاد}
نسل پنجم هم‌اکنون از تحلیل داده (\acro{NWDAF}) و هوش مصنوعی در بهینه‌سازی \acro{RAN} استفاده می‌کند. نسل ششم می‌خواهد اجزای یادگیرنده را بخشی از طراحی
کند: بازخورد وضعیت کانال، گیرنده‌ها و پیش‌بینی پرتو مبتنی بر هوش مصنوعی، و خودکارسازی شبکه با حلقه‌های بسته. پرسش‌های باز عبارت‌اند از استانداردسازی
مدل‌ها، تعمیم به شرایط دیده‌نشده، هزینهٔ انرژی خود هوش مصنوعی و توضیح‌پذیری.

% ---------------------------------------------------------------------
\section{مدل‌های ریاضی}

\subsection{افت مسیر فضای آزاد}\label{sec:fspl}
\paragraph{مدل} میان دو آنتن همسانگرد در فاصلهٔ $d$ و طول موج $\lambda = c/f$:
\begin{equation}\label{eq:fspl}
  L_{\mathrm{FS}} = 20\log_{10}\!\left(\frac{4\pi d f}{c}\right) \ \mathrm{dB}.
\end{equation}
با روزنه‌های مؤثر $A_t$ و $A_r$ معادلهٔ فریس به $P_r = P_t A_t A_r/(\lambda^2 d^2)$ تبدیل می‌شود که با $f^2$ \emph{افزایش} می‌یابد.
\paragraph{مثال} در ۱۰۰ متر: ۸۳٫۳ دسی‌بل در ۳٫۵ گیگاهرتز، ۱۰۱٫۴ دسی‌بل در ۲۸ گیگاهرتز و ۱۱۵٫۴ دسی‌بل در ۱۴۰ گیگاهرتز (\cref{fig:fspl}) -- ۳۲ دسی‌بل بیشتر در
۱۴۰ گیگاهرتز نسبت به ۳٫۵ گیگاهرتز. آرایه‌هایی با همان اندازهٔ فیزیکی در دو سو هر کدام $20\log_{10}(140/3.5) = 32$ دسی‌بل بهره می‌گیرند، پس بودجهٔ پیوند حتی
می‌تواند بهتر شود، به بهای پرتوهای بسیار باریک.

\begin{acnfigure}
  \acnfig{lecture-26/fspl}
  \caption{افت مسیر فضای آزاد (\cref{eq:fspl}).}
  \label{fig:fspl}
\end{acnfigure}

\subsection{پهنای باند تحت محدودیت توان}\label{sec:bw}
\paragraph{مدل} با توان دریافتی $P$ و چگالی نویز $N_0$:
\begin{equation}\label{eq:cap}
  C(B) = B \log_2\!\left(1 + \frac{P}{N_0 B}\right) \ \xrightarrow{B\to\infty}\ \frac{P}{N_0 \ln 2} .
\end{equation}
\paragraph{مثال} اگر \acro{SNR} روی ۱ گیگاهرتز ۲۰ دسی‌بل باشد: $C = 6.66$ گیگابیت بر ثانیه؛ با ۱۰ گیگاهرتز ۳۴٫۶؛ با ۱۰۰ گیگاهرتز ۱۰۰ گیگابیت بر ثانیه؛ حد ۱۴۴
گیگابیت بر ثانیه است (\cref{fig:bw}). پهنای باند بیشتر تنها تا وقتی بسیار کمک می‌کند که \acro{SNR} بالا بماند -- از این رو به بهرهٔ آرایه نیاز است.

\begin{acnfigure}
  \acnfig{lecture-26/bandwidth-capacity}
  \caption{ظرفیت بر حسب پهنای باند در توان دریافتی ثابت (\cref{eq:cap}).}
  \label{fig:bw}
\end{acnfigure}

\subsection{بهرهٔ \lr{RIS}}\label{sec:ris}
\paragraph{مدل} با $N$ عنصر که فازشان سیگنال‌های بازتابی را در گیرنده هم‌راستا می‌کند، دامنه‌ها به‌صورت همدوس جمع می‌شوند:
\begin{equation}\label{eq:ris}
  P_r \propto N^2 , \qquad G_{\mathrm{RIS}} = 20\log_{10} N \ \mathrm{dB}
\end{equation}
نسبت به یک عنصر.
\paragraph{مثال} $N = 1024$: ۶۰٫۲ دسی‌بل بیش از یک عنصر (فازهای تصادفی تنها $10\log_{10}N = 30.1$ دسی‌بل می‌دادند). این بهره باید بر حاصل‌ضرب دو افت مسیر غلبه کند.

\subsection{تفکیک‌پذیری حسگری}\label{sec:isac}
\paragraph{مدل} سیگنالی با پهنای باند $B$ دو هدف را در فاصله از هم جدا می‌کند اگر
\begin{equation}\label{eq:range}
  \Delta r = \frac{c}{2B} .
\end{equation}
\paragraph{مثال} $B = 1$ گیگاهرتز: ۱۵ سانتی‌متر؛ $B = 10$ گیگاهرتز: ۱٫۵ سانتی‌متر. بنابراین پهنای باندهای زیرتراهرتز حسگری در حد سانتی‌متر را ممکن می‌کنند.

\subsection{تأخیر ماهواره}\label{sec:leo}
\paragraph{مدل} برای ماهواره‌ای در ارتفاع $h$، شعاع زمین $R$ و زاویهٔ ارتفاع $\varepsilon$، فاصلهٔ مایل برابر است با
\begin{equation}\label{eq:slant}
  d(\varepsilon) = \sqrt{(R + h)^2 - R^2\cos^2\varepsilon} - R\sin\varepsilon ,
\end{equation}
و تأخیر یک‌طرفه $d/c$ است.
\paragraph{مثال} $h = 600$ کیلومتر: ۶۰۰ کیلومتر و ۲٫۰ میلی‌ثانیه در سمت‌الرأس، ۱۰۷۶ کیلومتر و ۳٫۶ میلی‌ثانیه در $30^\circ$، ۱۹۳۲ کیلومتر و ۶٫۴ میلی‌ثانیه در $10^\circ$
(\cref{fig:leo}). یک ماهوارهٔ زمین‌آهنگ (۳۵\,۷۸۶ کیلومتر) دست‌کم ۱۱۹ میلی‌ثانیه نیاز دارد.

\begin{acnfigure}
  \acnfig{lecture-26/leo-delay}
  \caption{تأخیر انتشار یک‌طرفه تا یک ماهوارهٔ \acro{LEO} (\cref{eq:slant}).}
  \label{fig:leo}
\end{acnfigure}

% ---------------------------------------------------------------------
\section{معماری شبکه و نمودارها}
یک سامانهٔ محتمل نسل ششم هستهٔ خدمت‌محور نسل پنجم (درس ۱۸) را نگه می‌دارد و آن را با کارکردهای حسگری و هوش مصنوعی گسترش می‌دهد، دسترسی ماهواره‌ای و
زمینی را زیر یک هسته یکپارچه می‌کند، و یک \acro{RAN} مستقر می‌کند که سلول‌های ماکرو در باند میانی بالایی، نقاط متراکم زیرتراهرتز، خوشه‌های بدون سلول و
پوشش با کمک \acro{RIS} را ترکیب می‌کند. \cref{fig:usage} دید خدمت را می‌دهد؛ \cref{fig:fspl,fig:bw,fig:leo} قیدهای فیزیکی را کمّی می‌کنند.

% ---------------------------------------------------------------------
\section{رویه‌های گام‌به‌گام}

\begin{procedurebox}{بررسی بودجهٔ یک پیوند زیرتراهرتز}\label{proc:budget}
\begin{steps}
  \item افت فضای آزاد را در بسامد حامل و فاصله محاسبه کنید (\cref{eq:fspl})؛ حاشیه‌های جذب و انسداد را بیفزایید.
  \item بهره‌های آرایهٔ فرستنده و گیرنده را از اندازهٔ آرایه‌ها بیفزایید.
  \item توان نویز $N_0B$ را همراه با عدد نویز گیرنده و سپس \acro{SNR} را محاسبه کنید.
  \item ظرفیت را محاسبه (\cref{eq:cap}) و با نرخ هدف مقایسه کنید.
\end{steps}
\end{procedurebox}

\begin{procedurebox}{داوری دربارهٔ یک ادعای نسل ششم}\label{proc:claim}
\begin{steps}
  \item آیا الزامی در یک سند \acro{ITU} یا \lr{3GPP} است یا یک هدف پژوهشی؟
  \item کدام کمیت فیزیکی آن را محدود می‌کند (پهنای باند، \acro{SNR}، تأخیر، انرژی)؟ آن را با یکی از مدل‌های این درس بررسی کنید.
  \item چه هزینه‌ای دارد (طیف، سایت‌ها، انرژی، پیچیدگی)؟
  \item به کدام سناریوی کاربرد \cref{fig:usage} خدمت می‌کند؟
\end{steps}
\end{procedurebox}

% ---------------------------------------------------------------------
\section{مثال‌های حل‌شده}

\begin{examplebox}{پیوند زیرتراهرتز کوتاه‌برد}\label{ex:thz}
\textbf{مسئله.} افت فضای آزاد را در ۳۰۰ گیگاهرتز روی ۱۰ متر محاسبه کنید.

\textbf{حل.} $20\log_{10}(4\pi\times10\times3\times10^{11}/3\times10^{8}) = 102.0$ دسی‌بل -- تقریباً برابر ۲۸ گیگاهرتز روی ۱۰۰ متر؛ به همین دلیل زیرتراهرتز برای
بردهای کوتاه مناسب است.
\end{examplebox}

\begin{examplebox}{پهنای باند بیشتر}\label{ex:bw}
\textbf{مسئله.} با توان \cref{sec:bw}، پهنای باند ۵۰ گیگاهرتز چه ظرفیتی می‌دهد؟

\textbf{حل.} $\mathrm{SNR} = 100/50 = 2$؛ $C = 50\log_2 3 = 79.2$ گیگابیت بر ثانیه: ۵۰ برابر پهنای باند حالت ۱ گیگاهرتز تنها ۱۲ برابر ظرفیت می‌دهد.
\end{examplebox}

\begin{examplebox}{رفت‌وبرگشت ماهواره}\label{ex:leo}
\textbf{مسئله.} یک \acro{UE} و یک دروازه هر دو یک ماهوارهٔ \acro{LEO} در ۶۰۰ کیلومتری را در $30^\circ$ می‌بینند. تأخیر انتشار رفت‌وبرگشت از طریق ماهواره (لولهٔ خمیده) چقدر است؟

\textbf{حل.} یک‌طرفه \acro{UE} تا ماهواره تا دروازه: $2\times3.58 = 7.2$ میلی‌ثانیه؛ رفت‌وبرگشت ۱۴٫۳ میلی‌ثانیه -- سازگار با خدمات تعاملی، برخلاف \acro{GEO} (بیش از
۴۷۷ میلی‌ثانیه).
\end{examplebox}

% ---------------------------------------------------------------------
\section{مقایسهٔ فناوری‌ها}
\begin{table}[htbp]
\centering\small
\caption{ویژگی‌های نسل پنجم و نامزدهای نسل ششم.}\label{tab:cmp}
\rowcolors{2}{acnmist}{white}
\begin{tabularx}{\linewidth}{@{}l Y Y@{}}
\toprule
\rowcolor{white} & \lr{5G (IMT-2020)} & \lr{6G (IMT-2030)}، نامزدها \\
\midrule
سناریوهای کاربرد & \lr{eMBB}، \lr{URLLC}، \lr{mMTC} & فراگیر، \acro{HRLLC}، انبوه، همه‌جایی، هوش مصنوعی، حسگری \\
طیف & زیر ۷ گیگاهرتز و ۲۴ تا ۷۱ گیگاهرتز & + باند میانی بالایی و زیرتراهرتز \\
معماری رادیو & سلول‌ها با \acro{MIMO} انبوه & سلول‌ها + بدون سلول، \acro{RIS} \\
حسگری & مکان‌یابی & حسگری و ارتباط یکپارچه \\
غیرزمینی & \acro{NTN} از انتشار ۱۷ & یکپارچه از ابتدا \\
هوش مصنوعی & تحلیل داده (\acro{NWDAF})، بهینه‌سازی \acro{RAN} & رابط هوایی و بهره‌برداری هوش‌مصنوعی‌بنیاد \\
وضعیت & مستقرشده & پژوهش و استانداردسازی اولیه \\
\bottomrule
\end{tabularx}
\end{table}

\paragraph{تکامل، نه انقلاب} نسل ششم نیز مانند همهٔ نسل‌های پیشین کارکردهای پایدار درس ۲۵ را دوباره به کار می‌گیرد و قابلیت‌های جدید را تا جایی که فیزیک و
اقتصاد اجازه دهند می‌افزاید.

% ---------------------------------------------------------------------
\section{نکات کلیدی}
\begin{takeawaybox}
\begin{itemize}
  \item \lr{IMT-2030} شش سناریوی کاربرد و چهار جنبهٔ فراگیر تعریف می‌کند؛ الزامات عددی بعداً می‌آیند.
  \item افت فضای آزاد با $20\log_{10} f$ رشد می‌کند، اما آرایه‌هایی با روزنهٔ ثابت آن را جبران می‌کنند: مسئله پرتو است، نه توان.
  \item پهنای باند تنها وقتی کمک می‌کند که \acro{SNR} بالا باشد: $C \to P/(N_0 \ln 2)$.
  \item بهرهٔ \acro{RIS} با $N^2$، تفکیک‌پذیری حسگری با $c/(2B)$ مقیاس می‌یابد؛ تأخیر \acro{LEO} چند میلی‌ثانیه است.
  \item هنگام مطالعه دربارهٔ نسل ششم، الزامات را از اهداف پژوهشی جدا کنید.
\end{itemize}
\end{takeawaybox}

% ---------------------------------------------------------------------
\section{تمرین‌ها}

\begin{exercisebox}{1}{سناریوها}\label{exr:scen}
سه سناریوی کاربرد \lr{IMT-2030} را که نسل پنجم را گسترش می‌دهند و سه سناریوی جدید را نام ببرید.
\end{exercisebox}

\begin{exercisebox}{1}{انسداد}\label{exr:block}
دو دلیل بیاورید که چرا پیوندهای زیرتراهرتز از پیوندهای ۳٫۵ گیگاهرتز به انسداد حساس‌ترند، و یک راه کاهش آن را بیان کنید.
\end{exercisebox}

\begin{exercisebox}{2}{تفکیک‌پذیری رادار}\label{exr:range}
تفکیک‌پذیری فاصلهٔ یک حامل \acro{NR} با پهنای ۴۰۰ مگاهرتز را که برای حسگری استفاده می‌شود محاسبه کنید.
\end{exercisebox}

\begin{exercisebox}{2}{اندازهٔ \lr{RIS}}\label{exr:ris}
یک \acro{RIS} شانزده‌عنصری بازتاب بسیار ضعیفی می‌دهد. برای ۲۰ دسی‌بل توان دریافتی بیشتر به چند عنصر نیاز است؟
\end{exercisebox}

\begin{exercisebox}{3}{\lr{GEO} در برابر \lr{LEO}}\label{exr:geo}
کمینهٔ تأخیر رفت‌وبرگشت \acro{UE} تا ماهواره تا دروازه را برای یک ماهوارهٔ \acro{GEO} در ۳۵\,۷۸۶ کیلومتری محاسبه و با \cref{ex:leo} مقایسه کنید. کدام خدمات شدنی‌اند؟
\end{exercisebox}

\begin{exercisebox}{3}{انرژی}\label{exr:energy}
ترافیک سالانه ۳۰٪ رشد می‌کند در حالی که بودجهٔ انرژی شبکه ثابت می‌ماند. انرژی به ازای بیت در هر سال و در ده سال چقدر باید بهبود یابد؟
\end{exercisebox}

\subsection*{پاسخ‌ها}
\addcontentsline{toc}{subsection}{پاسخ‌ها}

\begin{solutionbox}{\ref{exr:scen}}
گسترش‌یافته: ارتباط فراگیر، \acro{HRLLC}، ارتباط انبوه. جدید: اتصال همه‌جایی، هوش مصنوعی و ارتباط، حسگری و ارتباط یکپارچه.
\end{solutionbox}

\begin{solutionbox}{\ref{exr:block}}
طول موج کوتاه‌تر یعنی پراش کمتر پیرامون موانع و افت نفوذ بیشتر در مواد؛ پرتوهای باریک به‌آسانی مسدود می‌شوند. راه‌های کاهش: چنداتصالی به چند نقطهٔ
دسترسی، مسیرهای با کمک \acro{RIS}، تعویض سریع پرتو.
\end{solutionbox}

\begin{solutionbox}{\ref{exr:range}}
$\Delta r = 3\times10^{8}/(2\times4\times10^{8}) = 0.375$ متر.
\end{solutionbox}

\begin{solutionbox}{\ref{exr:ris}}
$(N/16)^2 = 100 \Rightarrow N = 160$ عنصر.
\end{solutionbox}

\begin{solutionbox}{\ref{exr:geo}}
هر گام دست‌کم ۳۵\,۷۸۶ کیلومتر (۱۱۹٫۳ میلی‌ثانیه): یک‌طرفه ۲۳۸٫۶ میلی‌ثانیه، رفت‌وبرگشت ۴۷۷ میلی‌ثانیه، حدود ۳۳ برابر مقدار \acro{LEO}. کیفیت صدا کاهش
می‌یابد و خدمات تعاملی یا کنترل بی‌درنگ شدنی نیستند؛ پخش همگانی و انتقال حجیم داده به‌خوبی کار می‌کنند.
\end{solutionbox}

\begin{solutionbox}{\ref{exr:energy}}
انرژی به ازای بیت باید هر سال با ضریب ۱٫۳ (۲۳٪) کاهش یابد، یعنی $1.3^{10} = 13.8$ برابر در ده سال -- از این رو بر حالت‌های خواب، سخت‌افزار کارآمد و
صرفه‌جویی انرژی مبتنی بر هوش مصنوعی تأکید می‌شود.
\end{solutionbox}

% ---------------------------------------------------------------------
\section{مطالعهٔ بیشتر}
چارچوب \lr{IMT-2030} توصیه‌نامهٔ \lr{ITU-R M.2160} است \cite{itu2023m2160}. سعد و همکاران \cite{saad2020vision} چشم‌اندازی از نسل ششم ارائه می‌دهند؛ راپاپورت و
همکاران \cite{rappaport2019above100} بسامدهای بالای ۱۰۰ گیگاهرتز را بررسی می‌کنند؛ انگو و همکاران \cite{ngo2017cellfree} \acro{MIMO} انبوه بدون سلول را معرفی
می‌کنند؛ وو و ژانگ \cite{wu2019irs} سطوح بازتابندهٔ هوشمند را تحلیل می‌کنند؛ لیو و همکاران \cite{liu2022isac} \acro{ISAC} را مرور می‌کنند؛ \lr{TR 38.811}
\cite{tr38811} \acro{NR} را برای شبکه‌های غیرزمینی مطالعه می‌کند.

\begin{latin}
\printbibliography[heading=subbibliography,title={References}]
\end{latin}
\end{document}
